In this section, we construct the estimation algorithm for the cost function and the control policy. 
This is based on taking the nonlinear optimal control problem and lifing it into a infinite dimensional linear representation. 
More specifically, we first formulate an infinite dimensional linear programming problem whose optimal solution is a cost function that corresponds to the observed data's occupation measure.
Next, we formulate an equivalent infinite dimensional linear programming problem which eleviates the problem of infinite time-horizon observations.
We then approximate the problem of recovering the cost function so that it is numerically tractable, and finally, based on a reconstrction of the cost function, we formulate an estimator for the control policy. Consistency of the estimators are proved in Sec.~\ref{sec:consistency}.


\subsection{The infinite dimensional IOC problem}
First note that if the running cost function $\ell$ and its corresponding value function $V$ satisfy
\begin{align*}
  \ell(x,a)-V(x)+\alpha\int_X V(y)Q(\mathrm dy|x,a) = 0,
\end{align*}
then any state-action trajectory $\{(\mfx_t,\mfa_t)\}_{t=1}^\infty$ is optimal to $\ell$.
For example, when the running cost $\ell(x,a)=\xi$ and the value function $V(x) = \frac{\xi}{1-\alpha}$, for some $\xi\geq0$ and all $(x,a)\in\mK$.
To avoid this trivial and uninformative solution to the IOC problem, we further make the following assumption.
\begin{assumption}\label{ass:IOC_normalization}
The ``true" running cost and the ``true'' value function $(\elltrue,\Vtrue)$ satisfy
\[
\int_{\mK}\left[\elltrue(x,a)-\Vtrue(x)+\alpha\int_X \Vtrue(y)Q(\mathrm dy|x,a)\right]\mathrm dx \mathrm da\ge 1.
\]
\end{assumption}
The above assumption rules out the case that the running cost is constant since it requires the admissible state-action set $\mK$ to at least contain a non-zero measure set whose components are not optimal.


Next, we construct an IOC algorithm that is based on the optimality conditions \eqref{eq:opt_cond_inf_dim} and the adjoint operators \eqref{eq:adjoints}. To this end, we first assume that $\mu^{\pitrue}$ is known (we approximate it with data in Sec.~\ref{sec:finite-dim-ioc}). Now, for every $(x,a)\in\mK$, we define the violation of the complementary slackness condition \eqref{eq:complementary_slackness} as $\psi:\mK\rightarrow\mR$ where
\begin{equation}
  \psi(x,a;\theta_\ell,V)\!:=\!\theta_\ell^T\varphi(x,a)\! +\alpha (q^*V)(x,a)\!-\!V(x).
\label{eq:complementary_slackness_violation}
\end{equation}
Then, we construct the IOC algorithm by minimizing the ``violation" of \eqref{eq:complementary_slackness}. Namely,
\begin{subequations}\label{eq:IOC_origin}
\begin{align}
  \min_{\substack{\theta_\ell\in\mR^{n_\ell}, V\in\mathscr{F}(X)\\\psi\in\mathscr{F}(\mK)}} 
  &\;\langle\psi,\mu^{\pitrue} \rangle\label{eq:violation_of_complementarity_slackness}\\
  \st  \qquad &\;\psi = \!\theta_\ell^T\varphi\! +\alpha (q^*V)\!-\!\proj^*V\label{eq:psi_def}\\
  &\;\psi(x,a)\!\ge\! 0,\forall (x,a)\in\mK,\label{eq:bellman_inequality_constraint}\\
  &\int_{\mK}\psi(x,a)\mathrm{d}x\mathrm{d}a \ge 1.\label{eq:regularization}\\
  &\|\theta_\ell\|_{1}\!\le\! \beta_\ell,\: \|\psi\|_{\infty,\mK}\!\le\! \beta_\psi,\label{eq:l_psi_norm_bounds}\\
  &\|V\|_{\infty,X}\!\le\! \beta_V.\label{eq:V_psi_norm_bounds}
\end{align}
\end{subequations}
Note that the norm bounds $\beta_\ell$, $\beta_V$ and $\beta_\psi$ in \eqref{eq:l_psi_norm_bounds} and \eqref{eq:V_psi_norm_bounds} can be chosen to be arbitrarily large. The norm constraint guarantees that the solutions to the optimization problem is always bounded, i.e., that we have a bounded estimator, which is important in the analysis later. Nevertheless, even without the constraint \eqref{eq:l_psi_norm_bounds} and \eqref{eq:V_psi_norm_bounds}, the following Proposition still holds.

\begin{proposition}\label{prop:well_posed_ioc_infinite_dim}
The problem \eqref{eq:IOC_origin} attains at least one optimal solution, and the optimal value is $0$. Moreover, for any optimal solution $(\theta_\ell^\star, V^\star, \psi^\star)$, $\pitrue$ is an optimal control policy to the optimal control problem \eqref{eq:forward_problem_obj_fun} with the running cost $\ell^\star:=\theta_\ell^{\star T}\varphi$. 
\end{proposition}

\begin{proof}
First, $\mu^{\pitrue}\in\mathcal{M}_+(\mK)$ and the constraint \eqref{eq:bellman_inequality_constraint} ensure that the objective function \eqref{eq:violation_of_complementarity_slackness} is bounded below by zero. 
Moreover, the optimal value zero can be attained since the ``true" $(\elltrue=\thetatrue_\ell^T\varphi, \Vtrue, \elltrue +\alpha (q^*\Vtrue) - \proj^*\Vtrue)$ satisfies the optimality condition \eqref{eq:opt_cond_inf_dim} and thus is both feasible to \eqref{eq:IOC_origin} and have objective function value zero. 

Next, let $(\ell^\star, V^\star, \psi^\star)$ be any optimal solution to \eqref{eq:IOC_origin}. Then $0 = \langle\psi^\star,\mu^{\pitrue} \rangle = \langle \ell^\star +\alpha (q^*V^\star) - \proj^*V^\star, \mu^{\pitrue}\rangle$, and since constraint \eqref{eq:bellman_inequality_constraint} is also satsified, $\pitrue$ satisfies the optimality conditions \eqref{eq:opt_cond_inf_dim}.
\end{proof}

In addition, for bounded $\elltrue$, its corresponding value function $\Vtrue$ is the limit of the well-known Value Iteration (VI), which is a contraction map on $ \mathscr{F}(X) $ with respect to supremum norm \citep[p. 52]{hernandez2012discrete}. 
Thus we have $\left\|\Vtrue\right\|_\infty< \infty$. 
Furthermore, by Assumption \ref{ass:transition_kernel} and \ref{ass:obj_func}, $\Vtrue$ is the fixed point of VI, this implies that for all $\eta_1:=(x_1,a_1),\eta_2:=(x_2,a_2)\in\mathbb{K}$, it holds that
\begin{align*}
&\left|\min_{a_1\in A}(\elltrue+\alpha q^* \Vtrue)(x_1,a_1)-\min_{a_2\in A}(\elltrue+\alpha q^* \Vtrue)(x_2,a_2)\right|\\
&\le \sup_{a\in A}\left|(\elltrue+\alpha q^* \Vtrue)(x_1,a)-(\elltrue+\alpha q^* \Vtrue)(x_2,a)\right|\\
&\le\sup_{a\in A}\left[\left|\elltrue(x_1,a)-\elltrue(x_2,a)\right|+\alpha\left|q^*\Vtrue(x_1,a)-q^*\Vtrue(x_2,a)\right|\right]\\
&\le \sup_{a\in A}\left[(L_\ell+\|\Vtrue\|_\infty L_Q)\|[x_1^T,a^T]^T-[x_2^T,a^T]^T\|_2\right]\\
&=\underbrace{(L_\ell+\|\Vtrue\|_\infty L_Q)}_{\beta_C}\|x_1-x_2\|_2
\end{align*}

and hence $\Vtrue\in \text{Lip}(\beta_C)$. 
Thus we can restrict ourselves to the set $V\in \text{Lip}(\beta_C)$ and $\psi\in \contfun(\mK)$, and the ``true" $\Vtrue$ and the corresponding $\bar{\psi}$ are still feasible when solving \eqref{eq:IOC_origin}. That is, we instead consider the problem
\begin{subequations}\label{eq:IOC_origin_cont_func_set}
  \begin{align}
    \min_{\substack{\theta_\ell\in\mR^{n_\ell}, V\in \contfun(X)\\\psi\in \contfun(\mK)}} &\;\;\langle\psi,\mu^{\pitrue} \rangle\\
    \label{eq:IOC_origin_cont_func_set_obj}
    \st \qquad &\;\;\eqref{eq:psi_def}-\eqref{eq:V_psi_norm_bounds},\\
    &V\in \text{Lip}(\beta_C).
    \label{eq:IOC_origin_V_lip_constraint}
  \end{align}
\end{subequations}
Analogously to the constraints \eqref{eq:l_psi_norm_bounds} and \eqref{eq:V_psi_norm_bounds}, we do not have to know the exact number of $\beta_C$ in practice since such bound on Lipschitz constant can be taken arbitrarily large.

Nevertheless, we do not have access to the measure $\mu^{\pitrue}$ in practice. One reason is that $\mu^{\pitrue}$ is defined as a sum of infinitely many occupation measures $\mu_t^{\pitrue}$ over time, yet we do not have infinite data length in practice. However, we can show that finite-time observations of the agent are sufficient to achieve the same effect as $\mu^{\pitrue}$.


\begin{proposition}[Finite-time IOC algorithm]\label{prop:finite_time_horizon_equivalent}
  Under Assumption \ref{ass:stationary_markov_policy}, for any $\nu\in\mathcal{M}_+(X)$, if $\support(\nu) =  \support(\proj\mu^{\pitrue})$, then solving \eqref{eq:IOC_origin_cont_func_set} is equivalent to solving
  \begin{subequations}\label{eq:IOC_approx1}
  \begin{align}
    \min_{\substack{\theta_\ell\in\mR^{n_\ell}, V\in \contfun(X)\\\psi\in \contfun(\mK)}} &\;\;\langle\psi, \nu \otimes\pitrue  \rangle\label{eq:IOC_approx1_obj_func}\\
    \st \qquad &\;\;\eqref{eq:psi_def}-\eqref{eq:V_psi_norm_bounds},\eqref{eq:IOC_origin_V_lip_constraint}, \label{eq:IOC_approx1_constraints}
  \end{align}
\end{subequations}
  where 
  \begin{align*}
    \langle f, \nu\otimes\pitrue\rangle := \int_X \left[ \int_A f(x,a)\pitrue(\mathrm{d}a|x) \right]\nu(\mathrm{d}x),\quad \forall f.
  \end{align*}
More specially, it holds that $\hat{\mu}^{\pitrue,N}:=\sum_{t=1}^N\mu_t^{\pitrue} = \left(\sum_{t=1}^N\proj\mu_t^{\pitrue}\right)\otimes\pitrue$, and under Assumption \ref{ass:init_value}, it holds that $\support\left(\sum_{t=1}^N\proj\mu_t^{\pitrue}\right)=\support( \proj\mu^{\pitrue}).$
\end{proposition}

\begin{pf}
  Since $\support(\nu) =\support(\proj\mu^{\pitrue})$, for any $f:X\rightarrow \mR_+$, it holds that
  \begin{align*}
    \int_Xf(x)(\proj\mu^{\pitrue})(\mathrm{d}x) = 0
    \Leftrightarrow\int_Xf(x)\nu(\mathrm{d}x) = 0.
  \end{align*}
  By Propostion \ref{prop:well_posed_ioc_infinite_dim}, the optimal value of \eqref{eq:IOC_origin} is zero, and hence the optimal value of \eqref{eq:IOC_origin_cont_func_set} is also zero. By following the same argument as that in the proof of Proposition \ref{prop:well_posed_ioc_infinite_dim}, it follows from $\nu\otimes\pitrue\in\mathcal{M}_+(\mK)$ and the constraint \eqref{eq:bellman_inequality_constraint} that the optimal value \eqref{eq:IOC_approx1} is bounded from below by zero and that the ``true" $(\elltrue, \Vtrue, \elltrue +\alpha (q^*\Vtrue) - \proj^*\Vtrue)$ attains the optimal value zero. Now suppose $(\ell^\star,V^\star, \psi^\star)$ is optimal to \eqref{eq:IOC_origin_cont_func_set}.  Thus we have
  \begin{align*}
    &(\ell^\star,V^\star, \psi^\star) \text{ minimizes \eqref{eq:IOC_origin},}\\
    \Rightarrow&\langle\psi^\star, \mu^{\pitrue}\rangle = \int_X \left[ \int_A \psi^\star(x,a)\pitrue(\mathrm{d}a|x) \right](\proj\mu^{\pitrue})(\mathrm{d}x) = 0, \\
    \Leftrightarrow&\langle\psi^\star,\nu\otimes\pitrue\rangle = \int_X \left[ \int_A \psi^\star(x,a)\pitrue(\mathrm{d}a|x) \right]\nu(\mathrm{d}x) = 0.
  \end{align*}
  Since \eqref{eq:IOC_origin_cont_func_set} and \eqref{eq:IOC_approx1} have the same constraints, $(\ell^\star,V^\star, \psi^\star)$ is also optimal to \eqref{eq:IOC_approx1}. Conversely, the same argument applies.
\end{pf}

\subsection{The finite dimension approximation of the IOC algorithm}\label{sec:finite-dim-ioc}
The above proposition implies that, instead of solving \eqref{eq:IOC_origin_cont_func_set}, we can spare ourselves from infinite data length and solve \eqref{eq:IOC_approx1} with $\nu\otimes\pitrue=\hat{\mu}^{\pitrue,N}$. Yet we still do not know the distribution $\hat{\mu}^{\pitrue,N}$ in practice, and hence we need to further approximate \eqref{eq:IOC_approx1} with $M$ I.I.D. samples $\{(x_t^i,a_t^i)\}_{t=1}^N$, $i=1:M$. In particular, we approximate the objective function \eqref{eq:IOC_approx1_obj_func} with empirical average, i.e.,
\begin{align}
  \langle\psi,\hat{\mu}^{\pitrue,N} \rangle\approx\frac{1}{M}\sum_{i=1}^M\sum_{t=1}^N\psi(\mfx_t^i,\mfa_t^i).
  \label{eq:emphirical_average}
\end{align}
Nevertheless, the approximated problem still optimizes over $V\in\contfun(X)$ and $\psi\in\contfun(\mathbb{K})$, which is infinite dimensional and thus difficult to solve. Moreover, the inequality constraint \eqref{eq:bellman_inequality_constraint} is challenging to handle since it needs to hold for all $(x,a)\in\mK$.
To address this issue, we approximate $\psi$ with a polynomial and use Putinar's Positivstellensatz \citep{putinar1993positive}\citep[Thm.~3.20]{laurent2009sums} to make the polynomial nonnegative.
More precisely, we approximate $\psi$ by a $2d_\psi$-th order polynomial on $\mK$ (without loss of generality, we consider only even-order polynomials to simplify the analysis to come), and represent the polynomial with Lagrange interpolation basis \citep{sauer1995multivariate}. Namely, for all $\eta := (x,a)\in\mK$, we have
\begin{align}\label{eq:polynomial_approx}
    \psi(\eta)\approx \hat{\psi}(\eta):= \underbrace{
    \begin{bmatrix}
    \hat\psi(\eta_{[1]}) &\cdots &\hat\psi(\eta_{[D_\psi]})
    \end{bmatrix}
    }_{\theta_\psi^T}\phi(\eta;\eta_{[\cdot]}),
\end{align}
where $D_\psi = C(n_x+n_u+2d_\psi, 2d_\psi)$ is the dimension of $2d_{\psi}$-th order polynomial space, $\eta_{[\cdot]}=\{\eta_{[1]},\ldots,\eta_{[D_\psi]}\}$ are the interpolation nodes, and the Lagrange polynomial basis $\phi(\eta;\eta_{[\cdot]})$ takes the form
\[
\phi(\eta;\eta_{[\cdot]}) = 
\begin{bmatrix}
\phi_1(\eta;\eta_{[\cdot]}) &\cdots &\phi_{D_\psi}(\eta;\eta_{[\cdot]})
\end{bmatrix}^T.
\]

\begin{remark}\label{remark:using-interpolation}
Note that any polynomial basis would work. Nevertheless, here we choose to work with Lagrange interpolation basis $\phi(\eta;\eta_{[\cdot]})$. The reason is that this leads to computation benefits; more details about this is found in Sec. \ref{sec:practical_implementation}.
\end{remark}

Next, to enforce the nonnegativity of $\hat{\psi}$, note that by Assumption \ref{ass:compact_support}, $\mK$ can be characterized by the semi-algebraic set $\{(x,a)|g_i(x,a)\ge 0,i=1:n_x+n_a\}$, where $g_i(x,a)= (x_i-x_{i,\min})(x_{i,\max}-x_i)$ and analogously for the bounds on $a$.
The quadratic module \citep[Eq.~3.13]{laurent2009sums} generated by $g_1,\ldots,g_{n_x+n_a}$ takes the form
\begin{align*}
\mathscr{Q}(g_{1:n_x+n_a})\!:=\!\{u_0\!+\!\!\!\sum_{i=1}^{n_x+n_a}\!\!u_i g_i \mid u_0,u_i\!\in\!\Sigma^2\}.
\end{align*}
Notably, $\hat{\psi}\in \mathscr{Q}(g_{1:n_x+n_a})$ means that $\hat{\psi}(\eta)\ge 0,\forall \eta\in\mK$.
Nevertheless, constraining $\hat{\psi}$ to be in $\mathscr{Q}(g_{1:n_x+n_a})$ might, a prior, be too restrictive. To see that it is in fact not the case, we first observe that $\mK$ is also Archimedean \citep[Def.~3.18]{laurent2009sums}. To see this, we consider the function 
\begin{align*}
f(\eta) = \sum_{i=1}^{n_x+n_a}g_i(x,a)=-\|\eta-\zeta\|_2^2+\mbox{constant},
\end{align*}
where $\zeta = [\zeta_1^x,\ldots,\zeta_{n_x}^x,\zeta_{1}^a,\ldots,\zeta_{n_a}^a]^T$ and $\zeta_i^x = \frac{x_{i,\min}+x_{i,\max}}{2}$,  $\zeta_i^a = \frac{a_{i,\min}+a_{i,\max}}{2}$. It holds that $f\in\mathscr{Q}(h_{1:n_x+n_a},g_{1:m})$. Moreover, the set $\{\eta|f(\eta)\ge 0\}$ is compact. Hence by \citep[Thm.~3.17, Def.~3.18]{laurent2009sums}, $\mK$ is Archimedean. 
By Putinar's Positivstellensatz \citep{putinar1993positive}\citep[Thm.~3.20]{laurent2009sums}, all polynomials that are positive on $\mK$ must belong to $\mathscr{Q}(g_{1:n_x+n_a})$.
In particular, this means that 
\begin{align}
\hat{\psi}(\eta)> 0, \; \forall \eta\in \mK \quad
\implies \quad \hat{\psi}\in\mathscr{Q}(g_{1:n_x+n_a}).\label{eq:epsilon_ineq1}
\end{align}
Therefore we can approximate \eqref{eq:bellman_inequality_constraint} to abitrary accuracy with \eqref{eq:epsilon_ineq1}. 
On the other hand, since \eqref{eq:epsilon_ineq1} says that $\hat{\psi}$ should belong to the quadratic module, it implies that the coefficients $\theta_\psi$ must belong to the cone $\left\{\theta\in\mR^{D_\psi}|\theta^T\phi\in \mathscr{Q}(g_{1:n_x+n_a})\right\}$, which, with a slight abuse of notation, we will denote as $\theta_\psi \in \mathscr{Q}(g_{1:n_x+n_a})$.

With the polynomial approximation \eqref{eq:polynomial_approx}, we can re-write the empirical average \eqref{eq:emphirical_average} as
\begin{align*}
    &\frac{1}{M}\sum_{i=1}^M\sum_{t=1}^N\hat\psi(\mfx_t^i,\mfa_t^i) = h^T\theta_\psi,
\end{align*}
where 
\begin{align*}
\bm h = \begin{bmatrix}
        \frac{1}{M}\sum_{i=1}^M\sum_{t=1}^N\phi(\mfx_t^i,\mfa_t^i;\eta_{[\cdot]})
    \end{bmatrix}.
\end{align*}
Similarly, we can further approximate the constraint \eqref{eq:regularization} as
\begin{align*}
    \int_\mK \hat\psi(\eta) \mathrm{d}\eta = \theta_\psi^T\underbrace{\begin{bmatrix}
        \int_\mK\phi(\eta;\eta_{[\cdot]})\mathrm{d}\eta
    \end{bmatrix}}_{d} \ge 1.
\end{align*}
In addition, the structure of the constraint \eqref{eq:psi_def} lends itself to an approximation by linear equations.
To see this, noticing that $V$ only depends on $x$, we hence approximate $V$ by $2d_V$-th order polynomial on $X$,
\begin{align*}
    V(x) \approx \hat{V}(x):=\theta_V^T r(x), \; \theta_V \in \mR^{D_V},
\end{align*}
where $D_V = C(n_x+2d_V, 2d_V)$, and $r(x)$ is a basis of $2d_V$-th order polynomials on $\mR^{n_x}$.
Consequently, the coordinate of $\proj^*V$ and $q^*V$ with respect to basis $\phi(\eta;\eta_{[\cdot]})$ can be calculated by their value at the interpolation points, i.e.,
\begin{align*}
    (\proj^*V)(\eta) &= \theta_V^T\underbrace{\begin{bmatrix}
        (\proj^*r)(\eta_{[1]}) 
        &\cdots 
        &(\proj^*r)(\eta_{[D_\psi]})
    \end{bmatrix}}_{=: G_1^T}\phi(\eta;\eta_{[\cdot]}),\\
    (q^*V)(\eta) &= \theta_V^T\underbrace{\begin{bmatrix}
        (q^*r)(\eta_{[1]}) 
        &\cdots 
        &(q^*r)(\eta_{[D_\psi]})
    \end{bmatrix}}_{=: G_2^T}\phi(\eta;\eta_{[\cdot]}).
\end{align*}
In a similar manner, we further do a coordinate change on $\ell$ to the same basis $\phi(\eta;\eta_{[\cdot]})$.
More specifically, by Assumption \ref{ass:obj_func}, 
\begin{align*}
\ell = \theta_\ell^T\varphi =\theta_\ell^T
\begin{bmatrix}
\varphi_1\\\vdots\\\varphi_{n_\ell}
\end{bmatrix}
\approx \theta_\ell^T
\underbrace{\begin{bmatrix}
\gamma_{\varphi_1}^T\\\vdots\\\gamma_{\varphi_{n_\ell}}^T
\end{bmatrix}}_{=: H^T}\phi,
\end{align*}
where $\gamma_{\varphi_i}^T\phi(\eta;\eta_{[\cdot]})$ is an approximation of the feature $\varphi_i$ in the Lagrange polynomial basis $\phi(\eta;\eta_{[\cdot]})$.
Therefore, the constraint \eqref{eq:psi_def} now can be approximated with the following linear constraints on the coefficients of the polynomial approximations
\begin{align}
    &\theta_\psi\!=\!\! \underbrace{\begin{bmatrix}H &\alpha G_2\!-\!G_1\end{bmatrix}}_{\Xi_\psi \in \mR^{D_\psi\times (n_\ell + D_V)}}\! \begin{bmatrix}
        \theta_\ell\\\theta_V
    \end{bmatrix}\!\!\Rightarrow\!\!
    \underbrace{\begin{bmatrix}
    I &-H &G_1\!\!-\!\!\alpha G_2
    \end{bmatrix}}_{=: G}\!
    \begin{bmatrix}
    \theta_\psi \\\theta_\ell \\ \theta_V
    \end{bmatrix}\!=\!0.\label{eq:psi_ell_V_equality_constraint}
\end{align}



As a brief summary, we can approximate \eqref{eq:IOC_origin} with the convex optimization problem
\begin{subequations}\label{eq:IOC_approx2}
\begin{align}
\min_{\theta_\psi\theta_\ell,\theta_V}&\;\;\bm h^T\theta_\psi\\
&\;\; \eqref{eq:psi_ell_V_equality_constraint}\\
&\;\; \theta_\psi\in\mathscr{Q}(g_{1:n_x+n_a}),\label{eq:theta_psi_constraint}\\
&\;\; d^T\theta_\psi\geq 1,\\
&\;\; \|\theta_\psi\|_{\infty}^2\le \beta_\psi^\prime,\label{eq:theta_psi_bounded}\\
&\;\;\|\theta_\ell\|_{\infty}^2\le \beta_\ell^\prime,\;\|\theta_V\|_{\infty}^2\le \beta_V^\prime.\label{eq:theta_V_ell_bounded}
\end{align}
\end{subequations}
Moreover, since all norms on finite dimensional vector spaces are equivalent \citep[Theorem~2.4-5]{kreyszig1991introductory}, the constraints \eqref{eq:l_psi_norm_bounds},\eqref{eq:V_psi_norm_bounds} and \eqref{eq:IOC_origin_V_lip_constraint} approximated in finite dimension is relaxed and simplized to \eqref{eq:theta_psi_bounded} and \eqref{eq:theta_V_ell_bounded}, where $\beta_\psi^\prime,\beta_\ell^\prime,\beta_V^\prime$ can be chosen arbitrarily large in practice.


\subsection{Expert policy reconstruction}\label{sec:policy-recovery}
Solving \eqref{eq:IOC_approx2} gives the esimates $(\hat{\ell}_{M,\mathbf{d}},\hat{V}_{M,\mathbf{d}})$ for the ``true" running cost and the value function $(\elltrue,\Vtrue)$, where $\mathbf{d}=(d_\psi,d_V)$. 
With the identified $\hat{\ell}_{M,\mathbf{d}}$, any off-the-shelf algorithm (e.g., model predictive control) that solves the forward optimal control problem can be used for control policy reconstruction.
As mentioned previously in Sec.~\ref{sec:problem_formulation}, such control policy reconstruction method should have better generalizability and robustness compared to behaviour cloning.
Nevertheless, in the worst case scenarios,
the reconstructed control policy may be different from the expert's policy.
This is because the ``true" cost function may not be identifiable and the identified cost function does not align with the ``true" one; or because the ``true" cost function may admit multiple optimal control policies.
Despite this, we can still get a reasonable control policy by taking the prior knowledge of the cost structure and the IOC result into account in the worst case scenarios.
To this end, note that by Assumption \ref{ass:stationary_markov_policy} and \eqref{eq:opt_cond_inf_dim}, it holds for the ``true" $(\elltrue,\Vtrue)$ and the expert control policy $\pitrue$ that
\begin{align*}
0&=\langle \bar{\psi}:=\elltrue+(\alpha q^*-\proj^*)\Vtrue,\mu^{\pitrue}\rangle+ \langle \| a - \pitrue(x) \|_2^2],\mu^{\pitrue}\rangle\\
&=\langle \bar{\psi},(\proj \mu^{\pitrue})\otimes\pitrue\rangle+\langle \| a - \pitrue(x) \|_2^2],\mu^{\pitrue}\rangle\\
&=\int_X \bar{\psi}(x, \pitrue(x)) (\proj\mu^{\pitrue})(\mathrm{d}x) + \langle\| a - \pitrue(x) \|_2^2,\mu^{\pitrue}\rangle.
\end{align*}
On the other hand, due to \eqref{eq:bellman_inequality_constraint}, for expert policy $\pitrue(x)$ and all deterministic stationary policies $\pi$ it holds that
\[
\int_X \bar{\psi}(x, \pi(x)) (\proj\mu^{\pitrue})(\mathrm{d}x) +\langle\| \pitrue(x) - \pi(x) \|_2^2,\mu^{\pitrue}\rangle\ge 0.
\]
Hence we parameterize $\pi$ by $\theta_\pi\in\Theta$, where $\Theta\subset\mR$ is compact and $\pi(x;\cdot):\Theta\rightarrow A$ is continuous for any $x$, for instance, a neural network $\pi(x;\theta_\pi)$ whose weights are $\theta_\pi$. And then we estimate $\pitrue$ by solving
\[
\min_{\theta_\pi\in\Theta} \; \langle\bar{\psi}(x, \pi(x;\theta_\pi)) +\| \pitrue(x) - \pi(x;\theta_\pi) \|_2^2,\proj\mu^{\pitrue}\rangle.
\]
Notably, by following an argument similar to the proof of Proposition \ref{prop:finite_time_horizon_equivalent}, it is equivalent to solving
\begin{align}
\min_{\theta_\pi} \; & \langle\bar{\psi}(x, \pi(x;\theta_\pi))+\| \pitrue(x) - \pi(x;\theta_\pi) \|_2^2,\proj\hat{\mu}_t^{\pitrue,N}\rangle
\label{eq:policy-approx-reconstruction-infinite}
\end{align}
Though we do not have access to neither the actual distribution $\hat{\mu}_t^{\pitrue,N}$, nor $\bar{\psi}$, we can still approximate the above problem. First, the integral can be approximated with the empirical mean of the I.I.D. observations of the optimal ``state-action pair" trajectories $\{(x_t^i,a_t^i)\}_{t=1}^N$, $i=1:M$. Next, we approximate $\bar{\psi}$ with the approximation $\hat{\psi}_{M,\mathbf{d}}$ that solves \eqref{eq:IOC_approx2}. Together, this gives the finite-dimensional optimization problem
\begin{align}
\min_{\theta_\pi\in \Theta} \! \frac{1}{M}\!\sum_{t=1}^N\!\sum_{i=1}^M\! \hat{\psi}_{M,\mathbf{d}}(\mfx_t^i, \pi(\mfx_t^i;\theta_\pi)) \!+\! \| \mfa_t^i \!-\! \pi(\mfx_t^i;\theta_\pi) \|_2^2.
\label{eq:policy_approx_reconstruction}
\end{align}

As a summary, we have constructed the infinite dimensional IOC algorithm \eqref{eq:IOC_approx1} and

a finite dimensional approximation \eqref{eq:IOC_approx2} that

can be solved numerically. 
More specifically, solving \eqref{eq:IOC_approx2} yields an approximation of the cost function as well as the value function that can lead to the ``true" policy.
As mentioned previously, with the identified cost function, one may apply any standard optimal control method to obtain a policy with good generalization capability.
In addition, solving \eqref{eq:policy_approx_reconstruction} provides another way of estimating the optimal control policy. The complete estimation procedure is described in Algorithm~\ref{alg:estimator}. 

\begin{algorithm}[!htpb]
    \caption{Estimation procedure for cost and policy}
    \label{alg:estimator}
    \begin{algorithmic}[1]
        \Require Data $\{(x_t^i,a_t^i)\}_{t=1}^N$, $i=1:M$, system dynamics $Q(\cdot|x,a)$, state and action space $\mK = X\times A$, polynomial degrees $(d_\psi,d_V)$, parametrization $\pi(\cdot;\theta_\pi)$.
        \Ensure $(\hat{\ell}_{M,\mathbf{d}},\hat{V}_{M,\mathbf{d}})$ and $\pi(\cdot;\hat{\theta}_\pi)$.
        \State Design interpolation the points $\eta_{[\cdot]}$ and the polynomial bases $r(x)$.
        \State Calculate the matrix $G$, defined in \eqref{eq:psi_ell_V_equality_constraint}, from system dynamics and sample moments vector $h$ from data.
        \State Obtain the estimate $(\hat{\ell}_{M,\mathbf{d}},\hat{V}_{M,\mathbf{d}})$ by solving \eqref{eq:IOC_approx2}.
        \State Obtain the estimate $\pi(\cdot;\hat{\theta}_\pi)$: either by solving \eqref{eq:policy_approx_reconstruction} to estimate the expert policy, or by applying any standard optimal control algorithm.
    \end{algorithmic}
\end{algorithm} 


